% Part: first-order-logic 
% Chapter: axiomatic-deduction 
% Section: proving-things-quant

\documentclass[../../../include/open-logic-section]{subfiles}

\begin{document}

\olfileid[pt]{fol}{axd}{prq}

\olsection{\usetoken{P}{derivation} com Quantificadores}

\begin{ex}
Apresentemos !!a{derivation} de $(\lforall[x][!A(x)] \land
\lforall[y][!B(y)]) \lif \lforall[x][(!A(x) \land !B(x))]$.

Primeiro, observe que
\begin{align*}
  (\lforall[x][!A(x)] \land \lforall[y][!B(y)]) & \lif \lforall[x][!A(x)]\\
  \intertext{é uma instância de \olref[prp]{ax:land1}, e}
  \lforall[x][!A(x)] & \lif !A(a) \\
  \intertext{de \olref[qua]{ax:q1}. Assim, por \olref[pro]{prop:chain}, sabemos que}
  (\lforall[x][!A(x)] \land \lforall[y][!B(y)]) & \lif !A(a) \\
  \intertext{é !!{derivable}. Analogamente, como}
  (\lforall[x][!A(x)] \land \lforall[y][!B(y)]) & \lif \lforall[y][!B(y)] \qquad\text{e}\\
    \lforall[y][!B(y)] & \lif !B(a)\\
    \intertext{são, respectivamente, instâncias de \olref[prp]{ax:land2} e \olref[qua]{ax:q1},}
    (\lforall[x][!A(x)] \land \lforall[y][!B(y)]) & \lif !B(a)\\
    \intertext{é derivável por \olref[pro]{prop:chain}. Usando uma instância adequada de \olref[prp]{ax:land3} e duas aplicações de~\MP, vemos que}
    (\lforall[x][!A(x)] \land \lforall[y][!B(y)]) & \lif (!A(a) \land !B(a))\\
    \intertext{é derivável. Agora podemos aplicar \QR{} para obter}
(\lforall[x][!A(x)] \land \lforall[y][!B(y)]) & \lif \lforall[x][(!A(x) \land !B(x))].
\end{align*}
\end{ex}

\end{document}
